% État de l'art — Vesuvius Challenge (base PaperDatabase « elder-scrolls » + docs du dépôt Elder-scrolls)
% Compilation : lualatex etat_de_l_art.tex (deux passes pour la table des matières)
\documentclass[11pt,a4paper]{article}

\usepackage{amsmath}
\usepackage{fontspec}
\setmainfont{DejaVu Serif}[Scale=0.92]
\setsansfont{DejaVu Sans}[Scale=0.92]
\setmonofont{DejaVu Sans Mono}[Scale=0.85]
\usepackage[french]{babel}
\usepackage[margin=2.3cm]{geometry}
\usepackage{booktabs,tabularx,array,longtable}
\usepackage{enumitem}
\usepackage{xcolor}
\usepackage[hidelinks]{hyperref}
\usepackage{microtype}

\setlist{itemsep=2pt,topsep=4pt}
\newcolumntype{L}{>{\raggedright\arraybackslash}X}
\renewcommand{\arraystretch}{1.15}

\newcommand{\lu}{\textsf{\small[lu]}}
\newcommand{\deduit}{\textsf{\small[déduit]}}
\newcommand{\commu}{\textsf{\small[communauté]}}
\newcommand{\nr}{\textit{n/r}}
\newcommand{\pherc}[1]{PHerc.~#1}

\title{État de l'art du Vesuvius Challenge\\[4pt]
\large Lecture virtuelle des rouleaux d'Herculanum en micro-CT :\\ ce qui a été publié, ce que la communauté a obtenu, ce qui reste ouvert}
\author{Projet Elder-scrolls — synthèse de la base \texttt{elder-scrolls}}
\date{15 novembre 2026}

\begin{document}
\maketitle

\begin{abstract}
\noindent Ce document fait le point sur la lecture virtuelle des papyrus carbonisés d'Herculanum
au 15 novembre 2026. Il croise deux sources. La première est la littérature scientifique
réunie dans la base \texttt{elder-scrolls} : 25 articles, dont 9 lus en entier et fichés.
La seconde est la veille communautaire du dépôt Elder-scrolls (Discord, Substack, dépôt
\texttt{villa}, août–novembre 2026). Trois constats ressortent. (1) Le signal d'encre carbonée
en CT d'absorption existe, mais il est faible et se joue aux échelles fines. Sur
profilométrie, il tombe sous Dice~0,70 au-delà d'environ 1~µm. Les modèles de la communauté
en retrouvent pourtant à 8–9~µm, en distillant des labels obtenus à 2,4~µm. (2) Le goulot
s'est déplacé de l'encre vers la segmentation : obtenir des surfaces grandes, contiguës et
justes dans les zones compressées. (3) La validation reste artisanale. La littérature publiée
ne propose ni test placebo ni calibration ; la communauté a mis au point des contrôles
pratiques (lecture inversée, ensembles, vérification de la surface). Aucun modèle de
super-résolution 9~→~2,4~µm n'a été entraîné ni publié : c'est une piste ouverte.
\end{abstract}

\tableofcontents

%======================================================================
\section{Périmètre, sources et niveau de preuve}
%======================================================================

\paragraph{Sources.}
\begin{itemize}
  \item \textbf{Base \texttt{elder-scrolls}} (PaperDatabase) : 25 articles. 9 ont une fiche
  complète (lecture intégrale, résultats chiffrés dans \texttt{benchmarks.csv}, 41 lignes) ;
  16 sont indexés sans fiche. Seuls les 9 articles fichés servent ici comme preuve ; les
  autres sont listés en annexe~\ref{sec:nonfiches}.
  \item \textbf{Dépôt Elder-scrolls} : \nolinkurl{docs/COMMUNITY_BRIEF.md} (veille Discord,
  Substack et \texttt{villa}, août–15~nov.~2026), \nolinkurl{docs/PRIZES.md},
  \nolinkurl{research/prior_work_super_resolution.md} (recherche d'antériorité sur la
  super-résolution, 15~nov.~2026).
\end{itemize}

\paragraph{Niveau de preuve.} Chaque affirmation porte une étiquette :
\lu{} pour un fait lu dans un article de la base (section citée) ;
\commu{} pour un résultat annoncé par un membre de la communauté, non relu par les pairs et
non reproduit ici ; \deduit{} pour une inférence de ce document. Plusieurs articles de la base
sont des preprints arXiv non relus par les pairs ; ils sont signalés. Une valeur non rapportée
s'écrit \nr.

\paragraph{Limites de la base.} \texttt{gaps.md} n'a pas encore été rempli ; les axes
\emph{super-résolution}, \emph{label-efficiency}, \emph{domain-transfer} et
\emph{validation-fp} ne comptent qu'un ou deux articles chacun ; aucun article sur le
pré-entraînement auto-supervisé 3D n'y figure. Sur ces sujets, l'état de l'art ci-dessous
repose presque entièrement sur la veille communautaire.

%======================================================================
\section{Le problème et la chaîne de traitement}
%======================================================================

Les rouleaux d'Herculanum, carbonisés par l'éruption de 79, ne peuvent plus être déroulés
physiquement. Leur encre, un noir de carbone, a presque la même densité que le papyrus
carbonisé : la tomographie par absorption ne la distingue pas à l'œil
\lu~\cite{mocella2015,parsons2023}. Deux verrous en découlent \lu~\cite{parsons2023}, §1 :
\begin{enumerate}
  \item \textbf{détecter l'encre} sans contraste d'absorption franc ;
  \item \textbf{segmenter et aplatir} un feuillet mince, compressé, froissé, enroulé sur des
  dizaines voire des centaines de spires.
\end{enumerate}

La chaîne de référence de l'équipe du Challenge en 2026 \commu{} est la suivante :
\begin{enumerate}
  \item prédiction de surface (nnU-Net \emph{m7}, U-Net 3D \emph{recto}) et champs
  \emph{Lasagna} (normales, phase d'enroulement, fibres) ;
  \item contraintes : patchs vérifiés, \emph{tracks}, fibres, ombilic (axe central) ;
  \item ajustement global de spirale (\texttt{fit\_spiral}, 30\,000 itérations) ;
  \item aplatissement (Lasagna, SLIM ou \emph{flatboi}) ;
  \item rendu \texttt{tifxyz} dans les deux sens de profondeur ;
  \item détection d'encre (Reader V3 par défaut).
\end{enumerate}
Cette chaîne descend directement du pipeline EduceLab (segmentation, texturation,
aplatissement dans Volume Cartographer, puis classifieur d'encre \emph{ink-ID})
\lu~\cite{parsons2023}, §3.1.

%======================================================================
\section{Imagerie et nature du signal d'encre}
%======================================================================

\subsection{Contraste de phase synchrotron (2015–2017)}

Les premières lettres lues dans un rouleau fermé l'ont été en tomographie par contraste de
phase, sans apprentissage \lu~\cite{mocella2015}. Le scan a été fait à l'ESRF (ID17), à
70~keV, avec une résolution spatiale mesurée de 100~µm. Sur \pherc{Paris.~4}, les auteurs
lisent des séquences de lettres et retrouvent presque tout l'alphabet grec, sauf X et
Ψ. L'hypothèse avancée est morphologique : l'encre forme un léger relief, « au moins une
centaine de microns » au-dessus des fibres. Segmentation et lecture sont manuelles ; aucune
métrique ni contrôle des faux positifs n'est rapporté, et les auteurs reconnaissent que
certains traits se confondent avec les fibres.

L'équipe CNR/ESRF étend l'approche à \pherc{375} et \pherc{495} \lu~\cite{bukreeva2016,bukreeva2017}.
Le scan est fait à 73~keV avec un pixel de 47~µm. Le déroulage procède par « scorzatura
virtuelle » : pelage par paquets de feuillets, aplatissement en deux temps, et les fibres
croisées servent de grille de contrôle de la distorsion. Les auteurs annoncent des segments
de texte allant jusqu'à 14 lignes et attribuent les rouleaux à Philodème (\emph{De
Rhetorica}, hypothèse paléographique). Ils listent eux-mêmes ce qui peut imiter le texte :
plis, trous, fissures, contamination, artefacts d'aplatissement \lu~\cite{bukreeva2017}. Ces
deux textes sont des preprints, sans paramètres ni code permettant de les reproduire.

\subsection{L'encre est visible en absorption, par sa morphologie (2019–2023)}

Parker et al.\ contestent l'idée que l'encre carbonée serait invisible en micro-CT
\lu~\cite{parker2019}. Sur un fantôme, l'encre apparaît comme une « bosse ». La couche
d'encre mesure 3–5~µm, plus 4–12~µm d'absorption dans le papyrus, soit 3–17~µm au total. Un
CNN 3D classe des sous-volumes pris sur la surface aplatie ; les labels viennent de
photographies recalées. Les auteurs en déduisent un voxel minimal entre 3 et 12~µm,
« probablement près de 3 ». Les résultats sont uniquement qualitatifs : une lettre
authentique (un sigma lunaire), et des faux positifs alignés sur les fibres.

EduceLab-Scrolls formalise ce travail en jeu de données et en protocole \lu~\cite{parsons2023}.
Le jeu dépasse 11~To : 16 scans de 6 fragments et 9 scans de rouleaux, de 2,23 à 23,9~µm.
Les labels viennent de photographies IR à 1000~nm recalées. L'évaluation est une validation
croisée à 8 plis sur des moitiés de fragments. Le papier introduit le F0.5 et la lecture par
un papyrologue : sur le fragment 34, 7 des 15 caractères de la vérité terrain sont retrouvés,
sans caractère inventé (§2.2.3). Les auteurs estiment visuellement le rappel à environ 50~\%
(confiance faible). Les plages de résolution jugées prometteuses vont de 3,24 à 7,91~µm
(§4.1.1). C'est un preprint ; ses Tables~1–2 sont illisibles dans l'extraction, donc leurs
valeurs ne sont pas citées ici.

\subsection{Combien de signal survit à la résolution des scans de prix ?}

L'étude la plus directe est celle d'Angelotti et al.\ (\emph{Sci.~Rep.}~2026)
\lu~\cite{angelotti2026topo}. Elle ne porte pas sur le CT mais sur la \textbf{profilométrie
optique} de fragments ouverts : la topographie de surface seule, échantillonnée à 0,34~µm.
Le modèle est un nnU-Net 2D, évalué en validation croisée à 5 plis sur 14 échantillons
(16 lettres) issus de 3 papyrus.

\begin{table}[h]
\centering\small
\begin{tabular}{@{}lccccccc@{}}
\toprule
Échantillonnage latéral (µm) & 0,34 & 0,68 & 1,02 & 1,36 & 2,04 & 3,40 & 10,88 \\
\midrule
Modèle entraîné à cette résolution & 0,890 & 0,817 & 0,754 & 0,586 & 0,516 & 0,483 & 0,467 \\
Modèle 0,34~µm, entrée dégradée & 0,899 & 0,881 & 0,758 & 0,634 & 0,494 & 0,029 & 0,007 \\
\bottomrule
\end{tabular}
\caption{Dice médian de la détection d'encre depuis la topographie, selon la résolution
(\cite{angelotti2026topo}, Table~1). En tenant un papyrus entier à l'écart, le Dice médian
tombe à 0,757, 0,817 et 0,475 selon le papyrus (Table~2).}
\label{tab:topo}
\end{table}

\noindent Ce qu'on peut en lire :
\begin{itemize}
  \item seules les résolutions $\le$~1,02~µm dépassent Dice~0,70 \lu ;
  \item à 0,68~µm, un modèle entraîné fin puis appliqué à une entrée dégradée fait mieux
  qu'un modèle entraîné à 0,68~µm (0,881 contre 0,817) \lu. C'est un argument pour un
  enseignant haute résolution, donc pour la distillation \deduit ;
  \item la transposition au CT n'est pas directe. La dégradation par moyenne de blocs ne
  reproduit ni la PSF d'un scanner, ni son bruit, ni le volume partiel ; et le feuillet visé
  en CT est interne au rouleau \deduit. Si le seuil d'environ 1~µm valait aussi en CT, le
  relief pur ne survivrait pas à 8–9~µm. Que la communauté lise quand même des lettres à
  9~µm suggère qu'un autre signal est en jeu : texture, densité, effet de volume
  \deduit.
\end{itemize}

\noindent \textbf{Ce qu'observe la communauté} \commu :
\begin{itemize}
  \item le modèle canonique 2,4~µm atteint environ 0,95 sur des données natives à 2,4~µm,
  mais environ 0,11 sur les mêmes données sous-échantillonnées à 9~µm (Chris Scheirer, score
  à l'échelle de la lettre) ;
  \item le flou en profondeur expliquerait l'essentiel de la perte (AndreasHad) ;
  \item la même région compressée de Paris~4 se lit bien mieux à 2,4~µm qu'à 7,91~µm, et
  1,1~µm donne une encre plus nette que 2,4~µm (page \emph{open problems} 2026) ;
  \item les données ESRF ramenées à 9,6~µm battent les anciennes données DLS à 7,91~µm
  (Forrest).
\end{itemize}

%======================================================================
\section{Détection d'encre par apprentissage}
%======================================================================

\subsection{Travaux publiés}

\begin{table}[h]
\centering\small
\begin{tabularx}{\linewidth}{@{}p{3.3cm}p{3.2cm}Lp{2.6cm}@{}}
\toprule
Travail & Données & Modèle et labels & Résultat \\
\midrule
ink-ID, Parker 2019 \cite{parker2019} & fantôme 12~µm ; 1 fragment & CNN 3D sur patchs de surface ; labels photo & qualitatif \\
ink-ID, EduceLab 2023 \cite{parsons2023} & 4 fragments, CV 8 plis & classifieur 3D de patchs ; labels IR recalés & 7/15 caractères (fr.~34) \\
vFFC, Quattrini 2023 \cite{quattrini2023} & fragments 3,2~µm ; fr.~1 en test & ResNet-34 3D + goulot vFFC + décodeur 2D ; labels IR & F0.5 = 0,47 (test) ; 0,70 / 0,60 (Kaggle public / privé, ensemble) \\
\bottomrule
\end{tabularx}
\caption{Détection d'encre dans la littérature de la base. Les trois sont des preprints, sauf Parker~2019 (PLoS ONE).}
\end{table}

\noindent Le goulot vFFC n'apporte presque rien : F0.5 de 0,47 contre 0,46 pour une
convolution 3D ou sans goulot, sur un seul fragment de test, sans répétition
\lu~\cite{quattrini2023}, Table~3. Les augmentations comptent davantage : F0.5 passe de 0,37
sans augmentation à 0,46 avec le groupe diédral (Table~1). Aucun de ces travaux n'est évalué
sur un rouleau entier ni à 8–9~µm.

\subsection{Modèles de la communauté (août–novembre 2026)}

\begin{longtable}{@{}p{2.4cm}p{3.2cm}p{9.2cm}@{}}
\toprule
Modèle & Auteur, date & Ce qui est rapporté \commu \\
\midrule
\endhead
\textbf{Reader V3} & KLAVIS, oct. & Défaut de la communauté. Lettres les plus claires sur 0841 et 1447 tenus à l'écart ; fausse encre sur 5,5~\% du papyrus vierge ; plancher de sortie à 0,5 (lissage des labels). Sur 0841, AUC 0,916 contre 0,772 pour \texttt{ink\_9um}. \\
Reader v2 & KLAVIS, 27 sept. & 9~µm natif, 6 rouleaux et fragments. AUC 0,834 (protocole 116~keV) ; fausse encre 10~\% (contre 58~\% Hecate, 43~\% \texttt{ink\_9um}). \\
\texttt{ink\_9um} & équipe (err), 9 août & Hybride 3D-2D entraîné sur du 2,4~µm ramené à 9,6~µm ; 14 checkpoints écartés jusqu'à 0,14 en précision équilibrée ; le défaut seed42 step-075000 est le pire sur 0139. \\
v8-in & YoussefNader, 29 sept. & Entraîné seulement sur du 9~µm au signal confirmé ; AUC zéro-shot 0,86 / 0,81 sur 1447. \\
Hecate & équipe, 15 sept. & 3D-2D distillé du modèle 3D Paris~4 ; bon rappel, beaucoup de faux positifs. \\
Rectina m1–m4 & evilaliv3 & Entraîné de zéro sur des labels audités de 8 rouleaux ; m2 proche de V3, traits épaissis. \\
\bottomrule
\end{longtable}

\paragraph{Recettes d'entraînement qui ont fait leurs preuves} \commu
\begin{itemize}
  \item \textbf{Labels à 2,4~µm, inférence à 9~µm.} La recette standard pour v2 et V3 :
  labels produits sur le 2,4~µm, données ramenées à environ 9,6~µm, inférence à 9~µm.
  \item \textbf{Pseudo-labels denses plutôt que labels manuels épars.} KLAVIS a utilisé les
  cartes de modèles 2,4~µm comme labels, soit environ 50~fois la surface manuelle. L'écart de
  score entre zone supervisée et zone non supervisée est passé d'environ 0,25 à 0,06. Les
  labels natifs denses de Nieuwlaar ont fait passer l'AUC tenue à l'écart de 0,915 à 0,955.
  \item \textbf{Les fragments ne transfèrent pas aux rouleaux.} L'entraînement sur photos IR
  de fragments ne se transpose pas aux scans de rouleaux (KLAVIS, environ 200 expériences).
  C'est une limite importante des travaux publiés, tous évalués sur fragments \deduit.
  \item \textbf{Labels « en anneau fermé»} (cœur d'encre, zone tampon ignorée, anneau
  négatif) ; IBN plutôt que BatchNorm. L'augmentation photométrique et l'adversarial de
  domaine (DANN) n'aident pas (jeff).
  \item \textbf{\texttt{scan\_transfer}} (TequilaYuen, PR \texttt{villa}~\#2512) : fait
  ressembler le 2,4~µm ramené à 9~µm à un vrai scan natif. Gain annoncé de +0,024 d'AUC
  (0139 tenu à l'écart, 3 graines), plus +0,011 avec TTA miroir et +0,018 avec
  l'auto-entraînement.
\end{itemize}

\paragraph{Débats ouverts} \commu
\begin{itemize}
  \item \textbf{3D ou 2,5D ?} L'équipe penche pour une encre détectée en 3D, sans
  aplatissement ; la plupart des soumissions gagnantes restent en 2,5D.
  \item \textbf{AUC ou lisibilité ?} Optimiser l'AUC épaissit les traits sans révéler de
  nouvelles lettres (brak, evilaliv3).
\end{itemize}

%======================================================================
\section{Segmentation de feuillets et déroulage}
%======================================================================

\subsection{Travaux publiés}

\begin{itemize}
  \item \textbf{Segmentation manuelle} (Barfod 2015, rouleau d'argent de Jerash, CT de
  laboratoire à 23~µm) \lu~\cite{barfod2015}. Les lignes sont dessinées à la main dans
  VGStudio MAX et dépliées par polylignes. Fort contraste métal/air : la méthode n'est pas
  transposable à l'encre carbonée.
  \item \textbf{Extraction locale itérative à garantie topologique} (Klenert 2023)
  \lu~\cite{klenert2023}. Fast marching depuis des graines, puis étiquetage en deux régions
  guidé par le tenseur de structure. La surface produite est garantie orientable, sur des
  couches serrées et repliées. Coût : 3~min sur un papyrus doublement plié, environ 0,2~s par
  graine ajoutée. Les volumes testés sont petits ($\sim 10^7$–$10^8$ voxels) et il n'y a pas
  de vérité terrain sur données réelles. La méthode pourrait se brancher sur une carte de
  probabilité apprise \deduit.
  \item \textbf{Ajustement difféomorphe de spirale} (Henderson, WACV~2026)
  \lu~\cite{henderson2025}. Approche descendante : une spirale d'Archimède est déformée par
  $T_\text{aff}\circ T_\text{flow}\circ T_\text{gap}$, ce qui garantit une surface sans
  auto-intersection. Elle est ajustée sur les prédictions nnU-Net de surface et de fibres. Sur
  la région centrale de Paris~4 (7,91~µm, 54~keV, environ 30 enroulements), elle bat
  ThaumatoAnakalyptor (tableau~\ref{tab:spiral}). Coût : 19~h sur une RTX~3090, dont 16~h de
  post-traitement. Code publié. Biais possible : les nnU-Net d'entrée ont été entraînés sur
  une partie de Paris~4 \deduit.
\end{itemize}

\begin{table}[h]
\centering\small
\begin{tabular}{@{}lccccc@{}}
\toprule
Méthode & WJF (\%) & MRWD & ChD & AD & Str \\
\midrule
Spirale difféomorphe (complet) & 3,20 & 7,64 & 5,29 & 0,0568 & 1,239 \\
\quad + SLIM & -- & -- & -- & 0,0933 & 1,031 \\
\quad sans n° d'enroulement relatifs & 2,77 & 25,67 & 7,72 & 0,0556 & 1,166 \\
ThaumatoAnakalyptor & -- & -- & 5,59 & 0,1567 & 1,027 \\
\bottomrule
\end{tabular}
\caption{Paris~4, région centrale, contre un maillage manuel ; plus bas = mieux
(\cite{henderson2025}, Table~1 ; colonnes reconstruites à l'extraction, confiance moyenne).
Le ChD est unidirectionnel : il ne pénalise pas une surface en excès.}
\label{tab:spiral}
\end{table}

\subsection{État de la pratique en 2026} \commu

\paragraph{Ce qui échoue encore.}
\begin{itemize}
  \item les régions compressées (« mushy »), les courbures serrées, les \emph{kolleseis} et
  les extrémités des rouleaux ;
  \item le raccord des enroulements entre fenêtres en $z$ ;
  \item les spirales qui se trompent sur le nombre d'enroulements ;
  \item l'accrochage cohérent au recto sur tout un rouleau.
\end{itemize}
Un maillage lisse étire l'encre d'un feuillet à l'autre, et les modèles d'encre entraînés
sur des surfaces \emph{m7} plus bosselées se dégradent sur des surfaces propres. Le
nnU-Net \emph{m7} rate les feuillets inclinés de plus d'environ 30° par rapport à son
orientation d'entraînement (Scheirer).

\paragraph{Méthodes émergentes.}
\begin{itemize}
  \item \textbf{Champs de nombre d'enroulement.} On cherche un champ scalaire dont les
  niveaux entiers sont les spires (JaimeLomRo, LimeGS, DJFiya, Wainer). La variante RGT-Est
  de Scheirer, sans apprentissage, réduirait d'environ 45~\% les fusions de feuillets.
  \item \textbf{Fibres.} Le solveur global d'enroulement par fibres de waldkauz (fusionné
  dans \texttt{villa}) traite un volume en environ 1~h. L'équipe y voit sa voie la plus
  prometteuse pour des contraintes à longue portée en zone difficile. S'y ajoutent
  l'extraction automatique de Lerognon et ScrollFiesta.
  \item \textbf{Représentations auto-supervisées.} \emph{dinovol} (DINO 3D, sean, 16~oct.)
  et la segmentation d'instances de feuillets HercuViT/HercUNet qui s'appuie dessus.
  \item \textbf{Pipelines de bout en bout.} \emph{ufsm} (Forrest) : feuillets lisibles sur
  Paris~3 et~4 seulement. Le pipeline de patchs de Will Stevens couvre au moins 33~\% de
  \pherc{1667}.
\end{itemize}

\paragraph{Diagnostic de l'équipe (sean, 14 nov.).} Les First Letters sont désormais limités
par l'obtention de segments grands, contigus et corrects, plus que par les modèles d'encre.
D'après l'équipe, de meilleurs modèles de surface viendront de labels dans les zones
difficiles, pas de nouvelles architectures.

%======================================================================
\section{Aplatissement et distorsion}
%======================================================================

\begin{itemize}
  \item Parker et al.\ (rouleau d'En-Gedi, preprint) comparent trois méthodes : LSCM, ABF et
  un modèle masses-ressorts (MM) \lu~\cite{parker2020}. Les métriques sont l'étirement
  $L^2$/$L^\infty$, l'erreur angulaire et l'erreur d'aire par triangle. ABF est le meilleur
  choix par défaut. LSCM déforme beaucoup les aires. MM s'étire plus et a replié un bord sur
  lui-même. Les valeurs chiffrées (annexe) ne sont pas dans le texte extrait.
  \item Henderson rapporte qu'un SLIM appliqué après l'ajustement de spirale ramène
  l'étirement de 1,239 à 1,031, au prix d'un défaut angulaire plus élevé (0,0568 →
  0,0933) \lu~\cite{henderson2025}.
  \item En pratique \commu{} : aplatissement Lasagna (\texttt{flatten\_spiral\_checkpoint.py}),
  SLIM ou \emph{flatboi}. Le critère de prix « paramétrisation isométrique » n'a pas de seuil
  chiffré : « ça doit avoir l'air juste à un humain ».
\end{itemize}
\textbf{Gap} : aucune mesure publiée ne relie la distorsion d'aplatissement à la qualité de
l'encre détectée ou à la lisibilité \lu~\cite{parker2020}, \deduit.

%======================================================================
\section{Résolution : super-résolution, distillation, transfert}
%======================================================================

\paragraph{Verdict de la recherche d'antériorité (15 nov.~2026).} Personne n'a entraîné ni
publié, résultats à l'appui, de modèle de super-résolution 9~→~2,4~µm (ou →~1,1~µm). L'idée a
été proposée plusieurs fois (2024 ; Will Stevens, 3~sept.~2026 ; Daviddemeij, 25~oct., dans
l'espace des prédictions). Forrest a publié le 23~oct. des recalages entre résolutions
« faits surtout pour la super-résolution », sans suite connue.

\paragraph{Ce qui fonctionne déjà} \commu
\begin{itemize}
  \item \textbf{Distillation 2,4~µm → 9~µm} : surface distillée de Ben Kyles (Progress Prize
  d'août), Reader v2/V3, étudiants de Forrest. La distillation de l'encre Paris~4 donne des
  faux positifs sur les bords à fort contraste.
  \item \textbf{Transfert de labels} et \texttt{scan\_transfer} (voir §4.2).
\end{itemize}

\paragraph{Littérature transférable.} Wang et al.\ (micro-CT de roches, ×4, 2D) obtiennent un
PSNR de 32 en L1 seul et de 29,7 avec un GAN \lu~\cite{wang2019}, §4.1. Le GAN donne une
texture plus plausible sans gain au pixel, et les auteurs eux-mêmes jugent légitime le
scepticisme envers la texture générée. C'est un avertissement direct pour l'encre, dont le
signal est justement une texture subtile \deduit.

\paragraph{Contraintes réglementaires} \commu. La preuve doit venir du volume éligible à
8–9~µm. Le 2,4~µm sert à l'entraînement et au repérage, et projeter un résultat 2,4~µm sur
le 9~µm est interdit. Le règlement du Grand Prix exclut « les données dérivées de scans à
plus haute résolution du volume soumis ». Aucune décision n'a été rendue sur un volume
super-résolu : il faut \textbf{demander aux organisateurs} avant de s'y engager.

\paragraph{Gap identifié.} Aucune comparaison contrôlée n'oppose « super-résoudre puis
détecter » à « entraîner le détecteur en basse résolution avec une supervision haute
résolution », ni en CT ni ailleurs. Les briques existent : recalages, lecteur de patchs
appariés \texttt{cross\_frame\_dataset.py}, dégradation synthétique, noyau
\texttt{scan\_transfer}, paires au même protocole (0139, 0814, 0841).

%======================================================================
\section{Validation et faux positifs}
%======================================================================

\paragraph{Littérature.} Aucun article fiché ne propose de test placebo, de lecture inversée
ou de calibration. L'évaluation va de l'inspection visuelle (2015–2019) à la validation
croisée sur moitiés de fragments (2023), avec un risque de fuite entre plis voisins
\deduit. Le seul score obtenu en tenant un papyrus entier à l'écart est celui de la
topographie : Dice médian de 0,475 à 0,817 selon le papyrus \lu~\cite{angelotti2026topo}.

\paragraph{Pratiques de la communauté} (ce qu'attendent les relecteurs) \commu
\begin{itemize}
  \item \textbf{Lecture dans les deux sens.} Une lecture inversée ne trouve pas d'encre réelle :
  0 fenêtre sur 48 dans le test d'injection de Sankur sur \pherc{800}. On garde le sens qui
  sépare, sans moyenner les deux à l'aveugle.
  \item \textbf{Vérifier d'abord que la surface est sur le papyrus.} Seules 81 des 201
  surfaces publiées sur les rouleaux éligibles montrent un feuillet. Un décalage de 4 à
  8~voxels coûte 0,07 à 0,22 d'AUC.
  \item \textbf{Ensembles.} Moyenner des modèles publics enlève l'essentiel du bruit, mais un
  ensemble de graines d'\texttt{ink\_9um} n'a aidé que sur 0 segment sur~6. Deux graines ne
  s'accordent que sur environ 16~\% de l'encre détectée (IoU).
  \item \textbf{Pièges.} Les textures « craquelées » sont rarement de l'encre. L'encre est
  plutôt \emph{plus claire} que le papyrus. Une fenêtre de modèle de 3–4~mm est acceptée si
  l'on montre qu'elle n'hallucine pas de formes de lettres.
  \item \textbf{Outils d'audit} : ink-or-hallucination, ink-placebo-check, le score tenu à
  l'écart ScrollGT (environ 0,75 d'AUC pour les modèles publiés), la vérité terrain caméra de
  KLAVIS sur 0343P.
\end{itemize}

%======================================================================
\section{Résultats obtenus à ce jour}
%======================================================================

\begin{table}[h]
\centering\small
\begin{tabularx}{\linewidth}{@{}p{1.8cm}p{2.3cm}p{3.2cm}L@{}}
\toprule
Date & Rouleau & Lauréat & Méthode \commu \\
\midrule
8 oct. & \pherc{343} & E.~Nieuwlaar & Spirale adaptée au scan 8,64~µm, ombilic propre ; encre = moyenne de 4 modèles 9~µm. \\
2 nov. & \pherc{306B} & L.~Warner, N.~Bodill & GrowPatch sur les prédictions m7 ; encre = produit TRACE v3B × fine-tune v8-in. \\
15 nov. & \pherc{211} & P.~Aranke & U-Net 3D résiduel sur 16 plans à 9~µm ; labels transférés d'un modèle 2~µm ; segment par spirale pleine hauteur, corrections manuelles et croissance guidée par DinoVol. \\
\bottomrule
\end{tabularx}
\caption{Prix First Letters attribués en 2026 (50\,000~\$ chacun). Il en reste 7, sur 19 rouleaux éligibles.}
\end{table}

\noindent Autres jalons \commu :
\begin{itemize}
  \item \pherc{1667} (Scroll~4) lu de bout en bout : un texte stoïcien, peut-être le plus
  ancien de la collection ;
  \item environ 140 colonnes de Paris~4 récupérées depuis le rescan à 2,4~µm, dont une
  correction interlinéaire. Le prix du titre reste ouvert et l'équipe a mis Paris~4 en pause
  le 26~oct. ;
  \item \pherc{139} identifié comme Philodème, \emph{Sur les dieux}, livre~8 ;
  \item 70~\% inférieurs de \pherc{172} déroulés.
\end{itemize}

%======================================================================
\section{Synthèse : verrous et pistes pour Elder-scrolls}
%======================================================================

\begin{table}[h]
\centering\small
\begin{tabularx}{\linewidth}{@{}p{3.2cm}LL@{}}
\toprule
Verrou & État de l'art & Ce qui manque \\
\midrule
Signal d'encre à 8–9~µm & Existe (Reader V3, AUC 0,92 sur 0841 tenu à l'écart \commu) ; borne profilométrique à environ 1~µm \lu & Mesure en CT réel de la part de signal perdue ; nature du signal exploité à 9~µm \\
Segmentation & Spirale + Lasagna ; champs d'enroulement et fibres en émergence & Zones compressées, kolleseis, raccords en $z$ ; labels en zones difficiles \\
Aplatissement & SLIM, ABF, Lasagna ; métriques Str/AD/$L^2$ & Lien entre distorsion et lisibilité \\
Résolution & Distillation 2,4→9~µm, \texttt{scan\_transfer} & Aucune super-résolution entraînée ; pas de comparaison SR vs distillation \\
Transfert entre rouleaux & Pseudo-labels denses multi-rouleaux ; fragments → rouleaux échoue & Explication des échecs (encres « différentes ») \\
Validation & Lecture inversée, ensembles, contrôle de surface & Calibration, incertitude, protocole standard ; littérature quasi vide \\
Auto-supervisé 3D & dinovol (communauté) & Aucun article dans la base \\
\bottomrule
\end{tabularx}
\caption{Verrous du domaine au 15 novembre 2026.}
\end{table}

\paragraph{État du dépôt Elder-scrolls.} Le dépôt compte 5 commits. La documentation de
contexte est complète (prix, veille, rouleaux, données, carte de \texttt{villa}, matériel). Il
contient des scripts de téléchargement et d'aperçu, et une recherche d'antériorité sur la
super-résolution. Aucun modèle n'a encore été entraîné : \texttt{src/scroll\_tools} est vide,
et l'environnement ML (\texttt{villa} sous WSL) n'est pas encore installé.

\paragraph{Pistes concrètes.}
\begin{enumerate}
  \item \textbf{Comparaison contrôlée SR contre distillation} sur des paires au même protocole
  (0139, 0814, 0841), avec des métriques de fidélité de l'encre plutôt que le PSNR. C'est le
  gap le plus net, et les briques existent. Demander d'abord aux organisateurs le statut d'un
  volume super-résolu.
  \item \textbf{Mesurer la perte de signal en CT}, en reproduisant le protocole d'Angelotti
  (tableau~\ref{tab:topo}) sur des paires recalées 2,4/9~µm au lieu de la profilométrie.
  \item \textbf{Viser les First Letters} avec la chaîne standard (spirale, Lasagna, Reader V3
  dans les deux sens), sur les 9 rouleaux qui n'ont toujours pas d'ombilic, en appliquant
  d'emblée les contrôles du §8.
\end{enumerate}

\paragraph{Lectures à compléter dans la base.}
\begin{itemize}
  \item Ficher Angelotti et al.\ 2026, \emph{Complete virtual unwrapping and reading of a
  rolled Herculaneum papyrus} (arXiv:2606.29085), probablement la référence centrale sur
  \pherc{1667} \deduit.
  \item Récupérer Seales et al.\ 2016 (En-Gedi, \emph{Science Advances}) et Brun et al.\ 2016
  (encre métallique, \emph{PNAS}), en libre accès mais bloqués au téléchargement automatique
  (liste \nolinkurl{A_TELECHARGER.md}).
  \item Ajouter des articles sur le pré-entraînement auto-supervisé 3D et sur l'hallucination
  en reconstruction tomographique (Bhadra et al., IEEE TMI 2021, cité dans la note sur la
  super-résolution).
\end{itemize}

%======================================================================
\appendix
\section{Articles de la base non encore fichés}
\label{sec:nonfiches}
%======================================================================
Ces articles sont indexés mais pas encore lus : rien n'est affirmé ici sur leur contenu.

\begin{small}
\begin{longtable}{@{}p{8.6cm}p{2.6cm}p{3.6cm}@{}}
\toprule
Titre & Venue, année & DOI / arXiv \\
\midrule
\endhead
Complete virtual unwrapping and reading of a rolled Herculaneum papyrus (Angelotti et al.) & arXiv, 2026 (preprint) & 2606.29085 \\
A model carbonized papyrus scroll opens a novel path to identifying readable scrolls of the Herculaneum Library (Seiler et al.) & PLoS ONE, 2026 & 10.1371/journal.pone.0353485 \\
Volumetric Texture Analysis for Unsupervised Ink Detection in Herculaneum Papyri (Korotetskyi, Haindl) & Zenodo, 2026 (rapport) & 10.5281/zenodo.20766168 \\
Xenon Contrast for Carbon-Ink Tomography (Hoffmann) & Zenodo, 2026 (rapport) & 10.5281/zenodo.23164128 \\
Pulsed thermographic analysis of Herculaneum papyri (Ceccarelli et al.) & Sci.~Rep., 2025 & 10.1038/s41598-025-19911-w \\
Deciphering scrolls with tomography: A training experiment (Foschiatti et al.) & arXiv, 2025 (preprint) & 2504.11485 \\
Layout of ancient Greek papyri through lead-drawn ruling lines revealed by Macro X-Ray Fluorescence Imaging (Romano et al.) & Sci.~Rep., 2023 & 10.1038/s41598-023-33242-8 \\
A computational platform for the virtual unfolding of Herculaneum Papyri (Stabile et al.) & Sci.~Rep., 2021 & 10.1038/s41598-020-80458-z \\
Unlocking history through automated virtual unfolding of sealed documents imaged by X-ray microtomography (Dambrogio et al.) & Nat.~Commun., 2021 & 10.1038/s41467-021-21326-w \\
Personalizing image enhancement for critical visual tasks: improved legibility of papyri (Atanasiu, Marthot-Santaniello) & IJDAR, 2021 & 10.1007/s10032-021-00386-0 \\
Browsing through sealed historical manuscripts by using 3-D CT with low-brilliance X-ray sources (Stromer et al.) & Sci.~Rep., 2018 & 10.1038/s41598-018-33685-4 \\
Virtual Recovery of Content from X-Ray Micro-Tomography Scans of Damaged Historic Scrolls (Rosin et al.) & Sci.~Rep., 2018 & 10.1038/s41598-018-29037-x \\
Tracking ink composition on Herculaneum papyrus scrolls: lead by X-ray techniques and Monte Carlo (Tack et al.) & Sci.~Rep., 2016 & 10.1038/srep20763 \\
\bottomrule
\end{longtable}
\end{small}

%======================================================================
\begin{thebibliography}{99}
\small
\bibitem{mocella2015} V.~Mocella, E.~Brun, C.~Ferrero, D.~Delattre. \emph{Revealing letters in rolled Herculaneum papyri by X-ray phase-contrast imaging}. Nature Communications 6:5895, 2015. \href{https://doi.org/10.1038/ncomms6895}{doi:10.1038/ncomms6895}.
\bibitem{barfod2015} G.~H.~Barfod, J.~M.~Larsen, A.~Lichtenberger, R.~Raja. \emph{Revealing text in a complexly rolled silver scroll from Jerash with computed tomography and advanced imaging software}. Scientific Reports, 2015. \href{https://doi.org/10.1038/srep17765}{doi:10.1038/srep17765}.
\bibitem{bukreeva2016} I.~Bukreeva et al. \emph{Enhanced X-ray-phase-contrast-tomography brings new clarity to the 2000-year-old `voice' of Epicurean philosopher Philodemus}. arXiv:1602.08071, 2016 (preprint ; version publiée : Sci.~Rep.~2016, à vérifier).
\bibitem{bukreeva2017} I.~Bukreeva, M.~Alessandrelli, V.~Formoso, G.~Ranocchia, A.~Cedola. \emph{Investigating Herculaneum papyri: An innovative 3D approach for the virtual unfolding of the rolls}. arXiv:1706.09883, 2017 (preprint).
\bibitem{parker2019} C.~S.~Parker, S.~Parsons, J.~Bandy, C.~Chapman, F.~Coppens, W.~B.~Seales. \emph{From invisibility to readability: Recovering the ink of Herculaneum}. PLoS ONE 14(5):e0215775, 2019. \href{https://doi.org/10.1371/journal.pone.0215775}{doi:10.1371/journal.pone.0215775}.
\bibitem{wang2019} Y.~D.~Wang, R.~T.~Armstrong, P.~Mostaghimi. \emph{Boosting Resolution and Recovering Texture of micro-CT Images with Deep Learning}. arXiv:1907.07131, 2019 (preprint ; version Water Resources Research à vérifier).
\bibitem{parker2020} C.~S.~Parker, W.~B.~Seales, P.~Shor. \emph{Quantitative Distortion Analysis of Flattening Applied to the Scroll from En-Gedi}. arXiv:2007.15551, 2020 (preprint).
\bibitem{parsons2023} S.~Parsons, C.~S.~Parker, C.~Chapman, M.~Hayashida, W.~B.~Seales. \emph{EduceLab-Scrolls: Verifiable Recovery of Text from Herculaneum Papyri using X-ray CT}. arXiv:2304.02084, 2023 (preprint).
\bibitem{quattrini2023} F.~Quattrini, V.~Pippi, S.~Cascianelli, R.~Cucchiara. \emph{Volumetric Fast Fourier Convolution for Detecting Ink on the Carbonized Herculaneum Papyri}. arXiv:2308.05070, 2023 (preprint).
\bibitem{klenert2023} N.~Klenert, V.~Lepper, D.~Baum. \emph{A Local Iterative Approach for the Extraction of 2D Manifolds from Strongly Curved and Folded Thin-Layer Structures}. arXiv:2308.07070, 2023 (preprint ; IEEE TVCG selon OpenAlex, non vérifié).
\bibitem{henderson2025} P.~Henderson. \emph{Virtually Unrolling the Herculaneum Papyri by Diffeomorphic Spiral Fitting}. WACV 2026 ; arXiv:2512.04927. Code : \url{https://github.com/pmh47/spiral-fitting}.
\bibitem{angelotti2026topo} G.~Angelotti, F.~Nicolardi, P.~Henderson, W.~B.~Seales. \emph{Ink Detection from Surface Topography of the Herculaneum Papyri}. Scientific Reports, 2026 ; arXiv:2603.27698. \href{https://doi.org/10.1038/s41598-026-58467-1}{doi:10.1038/s41598-026-58467-1}.
\bibitem{brief} Projet Elder-scrolls. \emph{Community brief: Vesuvius Challenge, Aug–Nov 2026} (\nolinkurl{docs/COMMUNITY_BRIEF.md}) et \emph{Prior work: super-resolution from $\sim$9~µm to 2.4~µm scans} (\nolinkurl{research/prior_work_super_resolution.md}), 15~nov.~2026. Sources primaires : Discord du Vesuvius Challenge, \url{https://scrollprize.substack.com}, \url{https://github.com/ScrollPrize/villa}.
\end{thebibliography}

\end{document}
